\documentclass[10pt]{article}
\usepackage[T1]{fontenc}
\usepackage{lmodern,amsmath,amssymb,amsthm,microtype}
\usepackage[margin=0.9in]{geometry}
\usepackage[colorlinks=true,urlcolor=blue,citecolor=blue]{hyperref}
\newtheorem{theorem}{Theorem}
\newtheorem{corollary}{Corollary}
\title{An Elementary Existence Correction and Finite Minimum Formula\\for OEIS A115038 and Its Companions}
\author{Qichao Wang\\{\small Hebei University of Technology}\\[0.3em]
Daoyu Dong\\{\small University of Electronic Science and Technology of China}}
\date{October 26, 2026}
\hypersetup{pdftitle={An Elementary Existence Correction for OEIS A115038},pdfauthor={Qichao Wang and Daoyu Dong}}
\begin{document}
\maketitle
\begin{abstract}
The existence conjecture recorded in A115038, A115039 and A115041 follows
from the elementary identity $((p-1)/2)^2+p=((p+1)/2)^2$ for odd primes
$p$, together with $5^2+2=3^3$. The least solution can be improved on
the square witness only by an odd prime power. This gives a finite exact
formula for every term and proves that all three companion sequences
are defined. The three entries encode the same existence question.
\end{abstract}

\section{Existence and the unique square witness}
For a prime $p$, let $X(p)$ be the least positive integer $x$ for which
$x^2+p=y^k$ with $y\ge2$, $k\ge2$. OEIS A115038 lists $X(p_n)$;
A115039 gives the largest exponent $k$ at that least $x$, and A115041
gives the corresponding base $y$ \cite{oeis}. All three entries currently
repeat the conjecture that a solution exists for every prime.

\begin{theorem}\label{thm:exist}
$X(2)=5$. For every odd prime $p$,
\[
 1\le X(p)\le (p-1)/2.
\]
The only positive solution of $x^2+p=y^2$ is
$x=(p-1)/2$, $y=(p+1)/2$.
\end{theorem}
\begin{proof}
For odd $p$, expand
\[
 \left(\frac{p-1}{2}\right)^2+p=
 \left(\frac{p+1}{2}\right)^2.
\]
Conversely $(y-x)(y+x)=p$ has the unique positive factor pair $(1,p)$,
which gives the asserted solution. For $p=2$, the values at $x=1,2,3,4$
are $3,6,11,18$, none a perfect power, while $5^2+2=27=3^3$.
\end{proof}

\newpage
\section{A finite minimum formula}
For an odd prime $p$, put $U=(p-1)/2$ and $B=(p+1)^2/4$.
For every odd prime $r\le\lfloor\log_2(B-1)\rfloor$, consider the finite
integer interval
\begin{equation}\label{eq:y}
 \left\lceil(p+1)^{1/r}\right\rceil
 \le y\le\left\lfloor(B-1)^{1/r}\right\rfloor.
\end{equation}
Let $\mathcal C(p)$ consist of $U$ together with the positive integers
$\sqrt{y^r-p}$ for pairs $(r,y)$ in these ranges for which $y^r-p$ is
a perfect square. Integer roots are taken exactly.

\begin{theorem}[Exact minimum]\label{thm:min}
For every odd prime $p$,
\begin{equation}\label{eq:min}
 X(p)=\min\mathcal C(p).
\end{equation}
Together with $X(2)=5$, this determines every term of A115038 by finite
root operations and finite minimization.
\end{theorem}
\begin{proof}
Every member of $\mathcal C(p)$ is an eligible positive $x$. If an
eligible $x$ is smaller than $U$, its value $V=x^2+p$ satisfies
$p<V<B$. Write $V=z^k$ with $k\ge2$, and choose a prime divisor $r$ of
$k$. If $r=2$ then $V$ is square, contradicting the uniqueness assertion
of Theorem~\ref{thm:exist}. Hence $r$ is odd, and
$V=(z^{k/r})^r$ gives a pair in \eqref{eq:y}. The lower bound follows
from $V\ge p+1$, the upper bound from $V\le B-1$, and $z^{k/r}\ge2$
implies $r\le\log_2(B-1)$. Thus every possible improvement appears.
\end{proof}

For the companion sequences set $V=X(p)^2+p$ and define
\[
 E(p)=\max\{k\in\{2,\ldots,\lfloor\log_2V\rfloor\}:
                \lfloor V^{1/k}\rfloor^k=V\},\qquad
 Y(p)=\lfloor V^{1/E(p)}\rfloor.
\]
These sets are finite and nonempty. Thus A115039$(n)=E(p_n)$ and
A115041$(n)=Y(p_n)$ exist at every positive index. Equivalently, $E(p)$
is the gcd of the exponents in the prime factorization of $V$.
For example $p=31$ gives $X(p)=1$, $V=32$, $E(p)=5$, $Y(p)=2$.

The existence identity is a standard difference-of-squares observation;
this note supplies a direct correction to the recorded conjecture and a
finite least-value reduction. The supplied verifier compares the formula
with an independent scan in $x$, using factorization to recognize perfect
powers and their maximal exponents.

\begin{thebibliography}{9}
\bibitem{oeis} C. Hilliard, OEIS A115038, A115039 and A115041,
\url{https://oeis.org/A115038}, \url{https://oeis.org/A115039},
\url{https://oeis.org/A115041}, accessed October 26, 2026.
\end{thebibliography}
\end{document}
